\scp{Summary of PMP Phonology}{01-07}
Before we proceed,
it is helpful for those unfamiliar with the conventions of
Austronesian scholarship for us to present an outline of PMP phonology.
Throughout this work we transcribe PMP with the symbols used by \citet{Blust2013},
as these are the symbols with the widest currency.
The only deviation is our use of ‹ə›
for the PMP mid-central vowel (traditionally transcribed ‹e›).

Certain symbols within this tradition of transcription are not standard IPA
and/or have unexpected phonetic values.
The symbols used for PMP reconstructions in this work
are given along with the phonetic values assigned to them by
\citet[553--591]{Blust2013} and \citet[31, 241]{Wolff2010}.
Although different scholars have suggested different
values for some of these consonants,
the system as a whole is broadly accepted.
Note in particular that the phonetic values of *z, *j, *r
and *R may be unexpected.

\begin{table}
	\centering\tcp{Proto-Malayo-Polynesian segments}{PMPSeg}
	\st{4.5pt}\begin{tabular}{llllllllllllllll}\lsptoprule
PMP	&	*p	&	*t	&	*c	&	*k	&	*q	&	*b	&	*d	&	*D	&	\mc{2}{l}{*z}	&	*j	&	*g	&	*m	&	*n	\\	\midrule
Blust	&	[p]	&	[t]	&	[ʧ]	&	[k]	&	[q]	&	[b]	&	[d]	&	[ɖ]	&	\mc{2}{l}{[ʤ]}	&	[gʲ]	&	[g]	&	[m]	&	[n]	\\	
Wolff 	&	[p]	&	[t]	&	---	&	[k]	&	[q]	&	[b]	&	[d]	&	---	&	\mc{2}{l}{[d̪ʲ]/[d̪ɟ]}	&	[g]	&	---	&	[m]	&	[n]	\\	\midrule \midrule
PMP	&	*ñ	&	*ŋ	&	*h	&	\mc{2}{l}{*s}	&	*l	&	*R	&	*r	&	*w	&	\mc{1}{l||}{*y}	&	*i	&	*ə	&	*a	&	*u	\\	\midrule
Blust 	&	[ɲ]	&	[ŋ]	&	[h]	&	\mc{2}{l}{[s]/[ʃ]}	&	[l]	&	[r]	&	[ɾ]	&	[w]	&	\mc{1}{l||}{[j]}	&	[i]	&	[ə]	&	[a]	&	[u]	\\	
Wolff 	&	---	&	[ŋ]	&	[h]	&	\mc{2}{l}{[t̪s̪]/[t̪θ]}	&	[l]	&	[ʁ]	&	---	&	[w]	&	\mc{1}{l||}{[j]}	&	[i]	&	[ə]	&	[a]	&	[u]	\\	
		\lspbottomrule
	\end{tabular}
\end{table}

\citet{Blust2013} only finds evidence for *c [ʧ] in western languages.
*D [ɖ] is only found word finally in some languages west of the Blust line.
Thus, neither of those proto-segments is directly relevant to our target region.
*ñ is not found word finally,
only initially and medially.
*j is not found word initially,
only medially and finally.
*y is extremely rare word initially.

\citet{Blust2013} proposes mid-vowels *e
and *o as a later development for his lower level
nodes of PCEMP and PCMP.
These are directly relevant to our target region,
and are discussed further in \srf{13-01-05}.

Several phonotactic combinations of proto-phonemes
listed in \trf{PMPSeg} often behave as a unit
and have separate reflexes from the individual proto-phonemes that comprise them,
thus requiring special attention in our target region.
These include:

\begin{itemize}
	\item Nasal-consonant clusters,
				abstractly symbolised as *NC.
				These often share the same point of articulation,
				such as *mp, *mb, *nt, *nd, *ns, *nz, *ŋk, *ŋg.
				They often behave as a unit,
				so it is not uncommon for the reflexes of *d to be different to those of *nd.
				Such as, for example,
				*d > \it{r} but *nd > \it{d}.
	\item In the middle of words,
				there are often two forms reconstructed (referred to as a doublet),
				with and without the nasal *(N)C, as in *upu and *umpu ‘grandparent,
				grandchild (reciprocal term)’.
				Both forms are needed to account for the data in the region,
				and are often found in neighbouring languages within the same subgroup.
	\item There is a stative prefix *ma- reconstructed for PMP.
				The presence of stative *ma- before roots beginning with *p or *b
				must be compared with corresponding *NC combinations.
				Thus, we pay special attention to *mp, *mb, *ma-p, and *ma-b.
				Sometimes these merge *mp/*mb/*ma-p/*ma-b > **mb
				in ways relevant for subgrouping at various levels,
				and sometimes different combinations of them reveal other notable patterns.
	\item Word-final vowel-glide sequences *-ay, *-aw, *-uy,
				and *-iw often have different reflexes
				from their individual proto-phonemes.
				They may behave as a unit,
				and not as a combination of a vowel and a consonant.
				Their reflexes are often relevant for subgrouping at different levels.
	\item Similar to the vowel-glide sequences above,
				word-final *-əj sometimes behaves as a unit,
				sometimes has a different conditioned reflex than schwa on its own,
				and is occasionally relevant for subgrouping at lower levels.
\end{itemize}

Most reconstructions in this book are from \citet{BlustTrussel2020},\footnote{
		After this book was drafted
		and the manuscript submitted for review
		an updated more interactive version of
		\citet{BlustTrussel2020} has become
		available as \citet{BlustSmith2023}.}
with two exceptions.
Firstly, for ‘water’ we usually use *wahiyəR from \citet[1027]{Wolff2010}
rather than \citeauthor{BlustTrussel2020}'s *wahiR.\footnote{
		\citet[1027]{Wolff2010} posits
		*wasiyeɣ (= *waSiyəR) at the level of
		Proto-Austronesian on the basis of MP cognates,
		as well as Atayal \it{qosiyaʔ} ‘water’
		and \ili{Taroko Sediq} \it{kesia}.}
With the possible exception of \tsc{\ili{Bima}-Lembata} languages (\crf{ch:BL}),
reflexes in languages of many subgroups in our target region
cannot derive their forms directly or in a straightforward way from *wahiR ‘water’,
so an alternate high level reconstruction with *y
and *ə, such as \citeauthor{Wolff2010} *wahiyəR,
or an intermediate form **wayəR is required.
Indeed, \citet[167]{Dempwolff1938} reconstructed *wayəR
(*vajeɣ in his original spelling).\footnote{
		While medial /h/ is required for some MP languages,
		such as Cebuano \it{wahig},
		many MP languages point to a form with medial *y and *ə.
		One example is \ili{Malay}.
		While the official Indonesian/Malaysian spelling of ‘water’ is now ‹air›,
		in many regions the pronunciation [aer] or [ayer] is dominant,
		and even reflected in older spellings such as ‹ajer› or ‹ayer›.
		These forms point to **wayəR rather than *wahiR.}

Secondly, we often make use of reconstructions below the
level of PMP which are marked with a double asterisk **.
Some of these reconstructions are PCEMP or PCMP
reconstructions from \citet{BlustTrussel2020},
while others are from other sources
and/or our own reconstructions.
Preliminary evidence for such lower-level reconstructions
is usually given in an accompanying note.
